\documentclass[10pt,a4paper]{amsart}
\usepackage[margin=19mm]{geometry}
\usepackage{amsmath,amssymb,mathtools}
\usepackage[scr=rsfs,cal=euler]{mathalfa}
\usepackage{xfrac}
\usepackage[unicode,hidelinks,bookmarks=false]{hyperref}
\newcommand{\ie}{i.e.\ }
\newcommand{\cf}{cf.\ }
\newcommand{\resp}{respectively\ }
\newcommand{\Subsection}{\subsection}
\providecommand{\nobreakdash}{\nobreak\hbox{-}}
\setlength{\emergencystretch}{2em}
\multlinegap0pt
\usepackage{bm}
\usepackage{bbm}
\usepackage{dsfont}
\usepackage{xspace}
\usepackage{cancel}
\usepackage{mathtools}
\usepackage{amsthm}
\makeatletter
\@ifundefined{thm}{\newtheorem{thm}{Theorem}}{}
\makeatother
\title[From brick manifolds to Grassmannians of bimodules]{From brick manifolds to Grassmannians of bimodules}
\author{Evgeny Feigin, Markus Reineke}
\date{Source: arXiv:2510.21319v1 (CC BY 4.0)}
\begin{document}
\maketitle

\noindent\emph{Source and licence.} From brick manifolds to Grassmannians of bimodules. Version arXiv:2510.21319v1 (CC BY 4.0), distributed under the Creative Commons Attribution 4.0 International licence, as recorded in the arXiv licence metadata for this exact version and shown on the abstract page https://arxiv.org/abs/2510.21319v1. Full rights notice: licenses/arxiv-2510.21319-CC-BY-4.0.txt.\par\smallskip

\noindent\textit{Editorial scope.} Counted unit: the Thm whose source label is \texttt{None} in the article (source0.tex lines 461--464, source label \texttt{None}), delivered together with the article's own complete original proof of it (source0.tex lines 465--510, one \texttt{proof} environment of 46 lines). No mathematics is written, repaired or re-derived: statement and proof are the article's own text line for line. Required context retained uncounted: none. Every definition, display and label of those lines is retained unchanged. Omitted from this package and not redistributed: 1--460, 511--913. Nothing inside a retained excerpt is omitted or altered: every retained body is delivered exactly as the article's lines are written, so no omission receipt of any kind accompanies this package. Citations: the retained excerpts carry no citation command at all, so no bibliography entry of the article is transcribed and no reference is redistributed. Reference adaptation: none was needed and none was made; every reference inside the delivered unit resolves inside it, so no unresolved reference or citation is produced. Printed slips kept literal: no printed word, symbol, hypothesis or number of the counted statement or of its proof was altered. Layout and class: the article's own class is \texttt{amsart} with packages including amssymb, amscd, hyperref, amsmath, amsthm, amsfonts, euscript, xy, tikz, tikz-cd; the wrapper is the lane's pinned \texttt{amsart} editorial wrapper at 10pt on A4, which restates the theorem-like environments the retained text needs and adds the article's own macro definitions that the retained text needs (none), with extra wrapper packages (bm, bbm, dsfont, xspace, cancel, mathtools). The delivered numbering is the unit's own. Assets: no retained span contains a figure environment, a graphics-inclusion command, a PDF-page inclusion, a table or a TikZ picture; no image file is copied into this package, the package declares no asset dependency and no image file is redistributed with it. Licence: version arXiv:2510.21319v1 of this article is distributed under the Creative Commons Attribution 4.0 International licence, as recorded in the arXiv licence metadata for this exact version and shown on the abstract page https://arxiv.org/abs/2510.21319v1. A full rights notice is delivered at licenses/arxiv-2510.21319-CC-BY-4.0.txt. Characters: every retained excerpt is pure ASCII, so the delivered engine, XeLaTeX with its default Latin Modern text font, reports no missing character.
\begin{thm}
There is a $G_{\bf d}$ equivariant emebedding $R_{\bf d}(Q)\subset X$ such that 
the image is dense in $X$.
\end{thm}

\begin{proof}
We consider the projective bimodule
$$N(V_*)=\bigoplus_{\alpha:i\rightarrow j}Ae_j\otimes V_i\otimes e_iA$$
and consider bimodule maps from $N(V_*)$ to $M(V_*)$ (inspired by the form of the standard projective resolution above).
%; the precise relation has to be explored!). 
We then have
\begin{multline*}
{\rm Hom}_{A\otimes A^{\rm op}}(N(V_*),M(V_*))\simeq\\
\bigoplus_{\alpha:i\rightarrow j}\bigoplus_{p\in Q_0}{\rm Hom}_{A\otimes A^{\rm op}}(Ae_j\otimes V_i\otimes e_iA,Ae_p\otimes V_p\otimes e_pA)\simeq\\
\bigoplus_{\alpha:i\rightarrow j}\bigoplus_{p\in Q_0}{\rm Hom}_k(V_i,{\rm Hom}_{A\otimes A^{\rm op}}(Ae_j\otimes e_iA,Ae_k\otimes e_kA)\otimes V_k)\simeq\\
\simeq\bigoplus_{\alpha:i\rightarrow j}{\rm Hom}_k(V_i, V_i\oplus V_j).
\end{multline*}
Here we use the fact that $${\rm Hom}_{A\otimes A^{\rm op}}(Ae_j\otimes e_iA,Ae_k\otimes e_kA)\simeq e_jAe_k\otimes e_kAe_i$$
is spanned by $\alpha\otimes e_i$ and $e_j\otimes \alpha$, making use of the assumption that there are no parallel paths in $Q$.

Similarly, we can compute
$${\rm End}_{A\otimes A^{\rm op}}(M(V_*))\simeq\bigoplus_{i\in Q_0}{\rm End}_k(V_i)$$
(even without the assumption on $Q$) and
$${\rm End}_{A\otimes A^{\rm op}}(N(V_*))\simeq\bigoplus_{\alpha:i\rightarrow j}{\rm End}_k(V_i)$$
(again using the assumption on $Q$).

Moreover, let us note that 
$${\rm\bf dim}\, M(V_*)-{\rm\bf dim}\, N(V_*)={\bf f}.$$
Namely, we have already proved that
$$M(V_*)_{i,j}=\bigoplus_{j\stackrel{\omega'}{\leadsto}p\stackrel{\omega}{\leadsto}i}V_p,$$
and similarly, we find
$$N(V_*)_{i,j}=\bigoplus_{j\stackrel{\omega'}{\leadsto}p\stackrel{\alpha}{\rightarrow}q\stackrel{\omega}{\leadsto}i}V_p.$$
The only summand of $M(V_*)_{i,j}$ not appearing as a summand of $N(V_*)_{i,j}$ is the summand $V_i$ corresponding to $j\stackrel{\omega'}{\leadsto}i\stackrel{e_i}{\leadsto}i$ in case there exists a path $\omega'$ from $j$ to $i$ (again using the assumption of $Q$ admitting no parallel paths). This proves the identity of dimension vectors.

Inside $X={\rm Gr}^{{\bf f}(V_*)}_{A\otimes A^{\rm op}}(M(V_*))$, we thus have the locally closed subset consisting of quotients whose kernel is isomorphic to $N(V_*)$. This subset is isomorphic to the quotient of the set
$${\rm Hom}^{\rm inj}_{A\otimes A^{\rm op}}(N(V_*),M(V_*))$$
of injective maps by the group
$${\rm Aut}_{A\otimes A^{\rm op}}(N(V_*)).$$
By the above calculation, this is nothing else than the quotient of
$$Z=\bigoplus_{\alpha:i\rightarrow j}{\rm Hom}_k(V_i, V_i\oplus V_j)$$
by the action of
$$H=\prod_{\alpha:i\rightarrow j}{\rm GL}(V_i).$$
On the open subset $Z^0$ of $Z$ consisting of tuples 
$$(F_\alpha=\left[F_\alpha^1\atop F_\alpha^2\right]:V_i\rightarrow V_i\oplus V_j)_{\alpha:i\rightarrow j}$$ 
where all $F^1_\alpha$ are invertible, we can use the action of $H$ to normalize all 
$F^1_\alpha$ to identity maps.  This shows that $Z^0/H\simeq R_{\bf d}(Q)$. 
The action of ${\rm Aut}_{A\otimes A^{\rm op}}(M(V_*))\simeq G_{\bf d}$ on $X$ 
then restricts to the standard change of base action of $G_{\bf d}$ on 
$$R_{\bf d}\simeq Z^0/H\subset Z/H\subset X,$$
yielding the claimed equivariant embedding in to $X$, which is of course a projective variety.
\end{proof}

\end{document}
